\begin{center}
\begin{minipage}{0.98\linewidth}
\centering
\textbf{先固定 $p=ab$，让两点的正参数相等}\par\medskip
\begin{tikzpicture}[font=\small,>=Stealth]

\begin{scope}[x=1.6cm,y=1.3cm]
\draw[->] (0,0)--(4,0) node[right] {$a$};
\draw[->] (0,0)--(0,3.8) node[above] {$b$};
\draw[CProof,very thick,domain=0.3:3.7,samples=100] plot (\x,{1/\x});
\draw[CExample,thick] (0,2)--(2,0);
\draw[dashed] (0,0)--(2.7,2.7) node[right] {$a=b$};
\fill[CProof] (1,1) circle (2.5pt);
\node[above right,fill=white] at (1.2,1.5) {$(\sqrt p,\sqrt p)$};
\node[right] at (2.9,0.6) {$ab=p$};
\node[below] at (2,0) {$2\sqrt p$};
\node[left] at (0,2) {$2\sqrt p$};
\end{scope}
\end{tikzpicture}\par\smallskip
图取 $p=1$。双曲线上 $a+b$ 在切线 $a+b=2\sqrt p$ 处最小，故 $(a+b)^2\ge4p$；再对 $p+k^2/p$ 取最小值，要求 $p=k$。
\end{minipage}
\end{center}
